% ============ 第1章 §1.6 Calderón-Zygmund 分解 (原书 pp.24-29) ============
\section{Calderón--Zygmund 分解}\label{sec:1.6}
\index{0symbolbad@bad part}

\index{0symbolgood@good part}

\index{E@二进制区间}








这里来解说 Calderón--Zygmund 分解. 正如其名所示, 这个分解是由 Calderón 与 Zygmund 构思的分解. 给定实数列 $a = \{a_{m}\}_{m=1}^{N}$ 时, 考虑截断操作. 也就是说, 在 $a_{m}$, $m = 1, 2, \ldots, N$ 中, 若绝对值大于给定的数 $K > 0$, 则将其替换为 $0$, 这样得到的数列记为 $c = \{c_{m}\}_{m=1}^{N}$. 对 $m = 1, 2, \ldots, N$, 令 $b_{m} = a_{m} - c_{m}$. 这样, 数列 $a$ 就借助 $b = \{b_{m}\}_{m=1}^{N}$ 与 $c$ 分解为 $a = b + c$. $c$ 在一定程度上继承了数列 $a$ 的性质, 但对于 $b$ 则什么也不知道. 若只要求 $c$ 的 $\ell^{\infty}$-范数不超过 $K$, 那么在 $a_{m}$, $m = 1, 2, \ldots, N$ 中, 当绝对值大于给定的数 $K > 0$ 时, 保持符号, 将绝对值替换为 $K$, 其余分量不作特别变更, 这样得到的数列 $d$ 也应该能达到目的. 不过, 若令 $a = d + e$, 则 $e$ 的结构无法再进一步研究, 这一方面与前面的 $b$ 相同. 于是, Calderón 与 Zygmund 提出了研究这个 $b$ 的结构的手段. 这就是 Calderón--Zygmund 分解.

这里以 $N$ 为自然数, 考虑下标从 $0$ 到 $2^{N}-1$ 的长度为 $2^{N}$ 的数列. 给定 $j = 0, 1, \ldots, N$ 与 $k = 0, 1, \ldots, 2^{N}-1$ 时, 定义
\begin{equation}\tag{1.14}\label{eq:1.14}
I(j, k) = \Z \cap [2^{j}k, 2^{j}(k+1)).
\end{equation}
称 $I(j,k)$ 为二进制区间. 另外, 当 $j = 1, 2, \ldots, N$ 时, 令
\[
I^{-}(j, k) = \Z \cap [2^{j}k, 2^{j}(k + 2^{-1})) = I(j-1, 2k),
\]
\[
I^{+}(j, k) = \Z \cap [2^{j}(k + 2^{-1}), 2^{j}(k+1)) = I(j-1, 2k+1).
\]
$I^{\pm}(j,k)$ 也是二进制区间.

\begin{example}\label{exa:1.12}
设 $N = 3$. 具体写出 $I(j,k)$, $j, k = 0, 1, 2$.

(1) $I(0,0) = \{0\}$, $I(0,1) = \{1\}$, $I(0,2) = \{2\}$,

(2) $I(1,0) = \{0,1\}$, $I(1,1) = \{2,3\}$, $I(1,2) = \{4,5\}$,

(3) $I(2,0) = \{0,1,2,3\}$, $I(2,1) = \{4,5,6,7\}$, $I(2,2) = \{8,9,10,11\}$.

进一步, 考察 $I^{\pm}(j,k)$, $j = 1, 2$, $k = 0, 1, 2$ 的情形.

(1) $I^{-}(2,0) = I(1,0)$, $I^{-}(2,1) = I(1,2)$, $I^{-}(2,2) = I(1,4)$,

(2) $I^{+}(2,0) = I(1,1)$, $I^{+}(2,1) = I(1,3)$, $I^{+}(2,2) = I(1,5)$.
\end{example}

下面的定理通过在数直线上画出二进制区间即可明白.

\begin{theorem}\label{thm:1.7}
当 $j, j' = 0, 1, \ldots, N$, $k, k' = 0, 1, \ldots, 2^{N}-1$ 时, $I(j,k)$ 与 $I(j',k')$ 互不相交, 或者一方包含另一方.
\end{theorem}

在下面的定理中, 要使用这个二进制区间的几何结构.

\begin{theorem}[Calderón--Zygmund 区间]\label{thm:1.8}
设 $\lambda > 0$ 与 $a = \{a_{j}\}_{j=0}^{2^{N}-1}$ 已给出. 考虑形如 $I(j,k)$ 的集合, 其中绝对值的平均, 即 $\# I(j,k)$ 个数 $\{|a_{l}|\}_{l \in I(j,k)}$ 的平均超过 $\lambda$. 更详细地说, 设下面出现的 $E$, $\mathcal{U}_{0}$ 的定义中的 $j$, $k$ 满足 $j = 0, 1, \ldots, N$, $k = 0, 1, \ldots$, 考虑
\[
\mathcal{U}_{0} = \left\{I(j,k) : \frac{1}{\# I(j,k)}\sum_{l \in I(j,k)} |a_{l}| > \lambda\right\}.
\]
对这个 $\mathcal{U}_{0}$, 令 $E = \bigcup I(j,k)$.

(1) 若 $j \in I(N,0) \setminus E$, 则 $|a_{j}| \leq \lambda$. 因此, $E \neq \emptyset$ 的充分必要条件是存在满足 $|a_{j}| > \lambda$ 的 $j$.

(2) 对 $\lambda > 0$, $\# E \leq \lambda^{-1}\|a\|_{\ell^{1}}$ 成立.
\end{theorem}

\begin{proof}
(1) 证其逆否命题. 若 $|a_{j}| > \lambda$, 则 $I_{0,j} = \Z \cap [2^{0}j, 2^{0}(j+1)) = \{j\} \in \mathcal{U}_{0}$. 因此, $j \in I_{0,j} \subset E$. 关于充分必要条件, 由于取绝对值不超过 $\lambda$ 的数列的各项的绝对值后再取平均仍不超过 $\lambda$, 所以是显然的.

(2) 利用定理~\ref{thm:1.7}, 从构成 $E$ 的属于 $\mathcal{U}_{0}$ 的区间中取一部分, 可将 $E$ 表示为互不相交的二进制区间族 $\mathcal{U} = \{I(j_{l}, k_{l})\}_{l=1}^{L}$ 的并. 对 $l = 1, 2, \ldots, L$,
\[
\#(I(j_{l}, k_{l})) \leq \lambda^{-1}\sum_{m \in I(j_{l}, k_{l})} |a_{m}|,
\]
所以把它们加起来即得结论.
\end{proof}

利用定理~\ref{thm:1.8} 的证明中出现的 $\mathcal{U}$, 来分解 $a$.

\begin{definition}[Calderón--Zygmund 分解]\label{def:1.4}
将定理~\ref{thm:1.8} 的证明中出现的 $E$ 表示为互不相交的二进制区间族 $\mathcal{U} = \{I(j_{l}, k_{l})\}_{l=1}^{L}$ 的并. 进一步, 构成 $E$ 的区间 $I$ 一定包含在 $I(j_{l}, k_{l})$, $l = 1, 2, \ldots, L$ 中的某一个中, 且可以认为真包含 $I(j_{l}, k_{l})$ 的二进制区间 $J$ 不属于 $\mathcal{U}_{0}$. 对属于 $\mathcal{U}$ 中区间 $I$ 的整数 $k$, 将 $a_{k}$ 替换为属于 $I$ 的数列 $\{a_{l}\}_{l \in I}$ 的平均 $\frac{1}{\# I}\sum_{l \in I} a_{l}$, 这样得到的数列记为 $g = \{g_{j}\}_{j=0}^{2^{N}-1}$. 对不属于 $E$ 的 $k \in I(N, 0)$, 令 $g_{k} = a_{k}$. 其中, 当 $\mathcal{U}$ 为空集时, 令 $g = a$. 另外, 令 $b = \{b_{j}\}_{j=0}^{2^{N}-1} = a - g$, $b_{I} = \{\chi_{I}(j)b_{j}\}_{j=0}^{2^{N}-1}$.

称 $g$ 为好的部分 (good part), $b$ 为坏的部分 (bad part).
\end{definition}

\begin{lemma}\label{lem:1.3}
若 $\lambda > 2^{-N}\|a\|_{\ell^{1}}$, 则 $\|g\|_{\ell^{\infty}} \leq 2\lambda$ 成立.
\end{lemma}

\begin{proof}
由定理~\ref{thm:1.8}(1), 对 $j \in E$ 证明 $|g_{j}| \leq 2\lambda$ 就足够了. 由假设,
$\frac{1}{2^{N}}\left|\sum_{j=0}^{2^{N}-1} a_{j}\right| \leq \frac{1}{2^{N}}\|a\|_{\ell^{1}} \leq \lambda$,
所以 $I(N,0) \notin \mathcal{U}$. 因此, 对每个 $I \in \mathcal{U}$, 可以考虑包含 $I$ 且长度为 $2$ 倍的二进制区间 $J$. 由 $J \notin \mathcal{U}_{0}$, $\frac{1}{2\# I}\sum_{j \in I} |a_{j}| \leq \frac{1}{\# J}\sum_{j \in J} |a_{j}| \leq \lambda$, 所以得到结论.
\end{proof}

\begin{lemma}\label{lem:1.4}
若 $b_{j} \neq 0$, 则存在二进制区间 $I$ 使得 $j \in I \in \mathcal{U}$.
\end{lemma}

\begin{proof}
这是因为 $g$ 与 $a$ 的差异出现在 $\bigcup_{I \in \mathcal{U}} I$ 上.
\end{proof}

\begin{example}\label{exa:1.13}
下面的表是对 $a_{0}$ 到 $a_{15}$ 为
\[
0, \ 4, \ 0, \ 4, \ 2, \ 0, \ 6, \ 0, \ 12, \ 0, \ 2, \ 0, \ 6, \ 6, \ 6, \ 0
\]
给出的数列 $a = \{a_{j}\}_{j=0}^{15}$ 进行如下计算得到的表.

\begin{itemize}
\item 在上段写下 16 项的数列.
\item 从左边数起, 对奇数号的数字, 取其与右邻数字的平均, 在该数字下面写下求得的平均.
\item 从左边数起, 对偶数号的数字, 该数字下面留空.
\item 当数字只剩下一个时结束.
\end{itemize}

\begin{center}
\begin{tabular}{cccccccccccccccc}
\hline
$0$ & $4$ & $0$ & $4$ & $2$ & $0$ & $6$ & $0$ & $12$ & $0$ & $2$ & $0$ & $6$ & $6$ & $6$ & $0$ \\
\hline
$2$ &  & $2$ &  & $1$ &  & $3$ &  & $6$ &  & $1$ &  & $6$ &  & $3$ &  \\
\hline
$2$ &  &  &  & $2$ &  &  &  & $3.5$ &  &  &  & $4.5$ &  &  &  \\
\hline
$2$ &  &  &  &  &  &  &  & $4$ &  &  &  &  &  &  &  \\
\hline
$3$ &  &  &  &  &  &  &  &  &  &  &  &  &  &  &  \\
\hline
\end{tabular}
\end{center}

(1) 设 $\lambda = \pi$. 基于定义~\ref{def:1.4} 求 $\mathcal{U}_{0}$, $\mathcal{U}$, $E$, $g$, $b$, $b_{I}$, $I \in \mathcal{U}$, 如下所示.
\begin{align*}
\mathcal{U}_{0} = \{ & \{1\}, \{3\}, \{6\}, \{8\}, \{12\}, \{13\}, \{14\}, \{8,9\}, \{12,13\}, \\
& \{8,9,10,11\}, \{12,13,14,15\}, \{8,9,10,11,12,13,14,15\}\},
\end{align*}
\[
\mathcal{U} = \{\{1\}, \{3\}, \{6\}, \{8,9,10,11,12,13,14,15\}\},
\]
\[
E = \{1, 3, 6, 8, 9, 10, 11, 12, 13, 14, 15\},
\]
\[
g: 0, 4, 0, 4, 2, 0, 6, 0, 4 = \frac{1}{8}\sum_{j=8}^{15} a_{j}, 4, 4, 4, 4, 4, 4, 4
\]
\[
b: 0, 0, 0, 0, 0, 0, 0, 0, 8, -4, -2, -4, 2, 2, 2, -4
\]
\[
b_{\{1\}} = b_{\{3\}} = b_{\{6\}} = 0, \quad b_{\{8,9,10,11,12,13,14,15\}} = b
\]

(2) 设 $\lambda = 4$. 基于定义~\ref{def:1.4} 求 $\mathcal{U}_{0}$, $\mathcal{U}$, $E$, $g$, $b$, $b_{I}$, $I \in \mathcal{U}$, 如下所示.
\[
\mathcal{U}_{0} = \{\{6\}, \{8\}, \{12\}, \{13\}, \{14\}, \{8,9\}, \{12,13\}, \{12,13,14,15\}\},
\]
\[
\mathcal{U} = \{\{6\}, \{8,9\}, \{12,13,14,15\}\}, \quad E = \{6, 8, 9, 12, 13, 14, 15\},
\]
\[
g: 0, 4, 0, 4, 2, 0, 6, 0, 6, 6, 2, 0, 4.5 = \frac{1}{4}\sum_{j=12}^{15} a_{j}, 4.5, 4.5, 4.5
\]
\[
b: 0, 0, 0, 0, 0, 0, 0, 0, 6, -6, 0, 0, 1.5, 1.5, 1.5, -4.5 = a_{15} - g_{15}
\]
\[
b_{\{6\}} = 0,
\]
\[
b_{\{8,9\}}: 0, 0, 0, 0, 0, 0, 0, 0, 6, -6, 0, 0, 0, 0, 0, 0
\]
\[
b_{\{12,13,14,15\}}: 0, 0, 0, 0, 0, 0, 0, 0, 0, 0, 0, 0, 1.5, 1.5, 1.5, -4.5
\]
\end{example}

\begin{exercise}\label{exer:1.21}
设 $\lambda = 2.5$. 对例~\ref{exa:1.13} 的数列, 基于定义~\ref{def:1.4}, 求 $\mathcal{U}_{0}$, $\mathcal{U}$, $E$, $g = \{g_{j}\}_{j=0}^{15}$, $b$, $b_{I}$, $I \in \mathcal{U}$.
\end{exercise}
